\input{preamble.tex}

\title{Отчет по лабораторной работе № 8. \\ Настройка сетевых сервисов. DHCP}
\author{Сущенко Алина Николаевна \\ НПИбд-01-23}
\date{2026}

\begin{document}

\maketitle
\newpage

\tableofcontents

\newpage
\section{Цель работы}
Приобретение практических навыков по настройке динамического распределения IP-адресов посредством протокола DHCP (Dynamic Host Configuration Protocol) в локальной сети.

\newpage
\section{Выполнение работы}

\subsection{Настройка DNS-сервера}
\begin{enumerate}
\item В логическую рабочую область проекта добавлен сервер \texttt{ansusenko-dns} и подключен к коммутатору \texttt{msk-ansusenko-sw-3} через порт Fa0/2 (Рис. \ref{01}):

\begin{center}
    \centering
    \includegraphics[width=\textwidth]{../images/01.png}
    \captionof{figure}{Логическая схема локальной сети с DNS-сервером.}
    \label{01}
\end{center}

\item В конфигурации сервера указан адрес шлюза 10.128.0.1, адрес сервера - 10.128.0.5 с маской 255.255.255.0.

\item Проверяем конфигурацию после настройки \texttt{ansusenko-dns} (Рис. \ref{02}):

\begin{center}
    \centering
    \includegraphics[width=\textwidth]{../images/02.png}
    \captionof{figure}{Настройка сервера.}
    \label{01}
\end{center}

\item Настроен сервис DNS: в конфигурации сервера выбрана служба DNS, активирована (флаг On) и добавлены ресурсные записи (Рис. \ref{03}):
  \begin{itemize}
    \item \texttt{dns.ansusenko.rudn.ru} $\rightarrow$ 10.128.0.5
    \item \texttt{file.ansusenko.rudn.ru} $\rightarrow$ 10.128.0.3
    \item \texttt{mail.ansusenko.rudn.ru} $\rightarrow$ 10.128.0.4
    \item \texttt{www.ansusenko.rudn.ru} $\rightarrow$ 10.128.0.2
  \end{itemize}

\begin{center}
    \centering
    \includegraphics[width=\textwidth]{../images/03.png}
    \captionof{figure}{Настройка DNS-сервера.}
    \label{03}
\end{center}

\end{enumerate}

\subsection{Настройка DHCP-сервиса на маршрутизаторе}
\begin{enumerate}
\item На маршрутизаторе \texttt{msk-ansusenko-gw-1} выполнена настройка DHCP-пулов для четырёх подсетей.(Рис. \ref{04}):

\begin{center}
    \centering
    \includegraphics[width=\textwidth]{../images/04.png}
    \captionof{figure}{Настройка для подсетей.}
    \label{05}
\end{center}

\item Просмотр информации о созданных пулах DHCP выполнен командой \texttt{show ip dhcp pool} (Рис. \ref{05}):

\begin{center}
    \centering
    \includegraphics[width=\textwidth]{../images/05.png}
    \captionof{figure}{Информация о пулах DHCP.}
    \label{05}
\end{center}

\item Проверка привязок выданных адресов выполнена командой \texttt{show ip dhcp binding}. На момент проверки привязки отсутствуют, так как оконечные устройства ещё не запросили адреса. (Рис. \ref{06}):

\begin{center}
    \centering
    \includegraphics[width=\textwidth]{../images/06.png}
    \captionof{figure}{Информация о привязках DHCP.}
    \label{06}
\end{center}

\end{enumerate}

\subsection{Настройка оконечных устройств на получение адресов по DHCP}
\begin{enumerate}
\item На оконечных устройствах в настройках интерфейса изменён способ получения IP-адреса с \texttt{Static} на \texttt{DHCP}. Устройство \texttt{dk-ansusenko-01} получило IP-адрес 10.128.3.30, маску подсети 255.255.255.0, шлюз по умолчанию 10.128.3.1 и DNS-сервер 10.128.0.5. (Рис. \ref{07}):

\begin{center}
    \centering
    \includegraphics[width=\textwidth]{../images/07.png}
    \captionof{figure}{Настройка интерфейса устройства \texttt{dk-ansusenko-01} на получение адреса по DHCP.}
    \label{07}
\end{center}

\item Аналогичная настройка выполнена для остальных оконечных устройств в соответствующих подсетях.

\end{enumerate}

\subsection{Проверка доступности устройств из разных подсетей}
\begin{enumerate}
\item Выполнена проверка доступности между устройствами из разных подсетей с помощью утилиты \texttt{ping} (Рис. \ref{08}):

\begin{center}
    \centering
    \includegraphics[width=\textwidth]{../images/08.png}
    \captionof{figure}{Результат проверки доступности между устройствами из разных подсетей.}
    \label{08}
\end{center}

Пинг с устройства \texttt{dk-ansusenko-01} (10.128.3.30) на устройство из подсети departments (10.128.4.30) выполнен успешно, что подтверждает корректную работу маршрутизации и DHCP.

\item Так же пропингуем устройства из разных подсетей. \texttt{ping} (Рис. \ref{09}):

\begin{center}
    \centering
    \includegraphics[width=\textwidth]{../images/09.png}
    \captionof{figure}{Пинг устройства разных подсетей}
    \label{09}
\end{center}

\item В режиме симуляции изучим, каким образом происходит запрос адреса по протоколу DHCP.(Рис. \ref{10}): 

\begin{center}
    \centering
    \includegraphics[width=\textwidth]{../images/10.png}
    \captionof{figure}{Движение пакетов в режиме симуляции}
    \label{10}
\end{center}

\end{enumerate}

\section{Вывод}
В результате выполнения лабораторной работы приобретены практические навыки по настройке динамического распределения IP-адресов с использованием протокола DHCP. Настроен DNS-сервер, на маршрутизаторе созданы DHCP-пулы для четырёх подсетей, оконечные устройства успешно получили адреса по DHCP. Проверка доступности между устройствами из разных подсетей подтвердила корректность выполненных настроек.

\end{document}
